\documentclass[12pt]{article}
\usepackage[utf8]{inputenc}
\usepackage[T1]{fontenc}
\usepackage{hyperref}
\title{Tahltan Language Resources for K-12 Curriculum Development}
\author{Compiled by Claude Sonnet 5 (claude-sonnet-5) for Andr\'as Kornai}
\date{}
\begin{document}
\maketitle

\section{Introduction}
Tahltan (ISO 639-3 tht) is an endangered Athabaskan language of northern British Columbia, Canada. Searches found no formal K-12 curriculum. They did find a dictionary project run by the Tahltan Central Government (TCG), an older children's dictionary, and early-childhood language programmes. This report covers only what appeared in search results.

\section{Primary Educational Resources}
\begin{itemize}
\item Tahltan Children's Illustrated Dictionary (1994), by Colin Carter with elders. It is a Tahltan--English dictionary, openly available on FirstVoices: \url{https://wiki.ubc.ca/Documentation:RelLex/Tahltan_Children%27s_Illustrated_Dictionary}
\item Language nests for young children, the first opened in 2014 at Iskut. The sources give only brief mentions: \url{https://languageconservancy.org/a-dictionary-project-is-underway-to-revive-and-make-language-more-accessible-for-tahltans-terrace-standard/}
\end{itemize}

\section{Linguistic and Reference Materials}
\begin{itemize}
\item A manuscript grammar and dictionary of the language of the Tahltans on the Stikine River, 1957--1962, held by the University of Chicago Library and partly open access: \url{https://wiki.ubc.ca/Documentation:RelLex/A_Grammar_and_Dictionary_of_the_Language_Spoken_by_the_Tahltans_on_the_Stikine_River,_British_Columbia}
\item General overview: \url{https://en.wikipedia.org/wiki/Tahltan_language}
\end{itemize}

\section{Religious and Community Resources}
No scripture or religious materials turned up in the searches. The Tahltan Central Government is the community body behind the dictionary project: \url{https://languageconservancy.org/taltan-dictionary-project-aims-to-preserve-endangered-indigenous-language-cftk-tv/}

\section{Digital and Audio Resources}
\begin{itemize}
\item Tahltan FirstVoices entry: \url{https://wiki.ubc.ca/Documentation:RelLex/T%C4%81%C5%82t%C4%81n_FirstVoices}
\item TCG multimodal dictionary project, with audio, about 10,000 words and a planned offline app, supported by The Language Conservancy. The most recent coverage found dates from about 2022, so its current status is unconfirmed: \url{https://languageconservancy.org/?p=202832}
\end{itemize}

\section{Curriculum Development Context}
Only about 45 speakers remain, mostly elders. Existing material is a children's dictionary, a manuscript grammar and a dictionary project in progress. Textbooks, graded readers and a school curriculum were not found. A K-12 programme would need to build these with TCG and the fluent speakers.

\end{document}